\chapter{Thinking Functionally: From Imperative Loops through OCaml and Haskell to Rust}
\label{ch07:functional-thinking}

In Part~II we opened the hood on Rust's memory and type machinery: single ownership
and borrowing in Chapter~3, region lifetimes in Chapter~4, algebraic \texttt{struct}
and \texttt{enum} layouts in Chapter~5, and Chapter~6's separation of flat data from
behavior through traits such as \texttt{FeePolicy}, static monomorphization, and
16-byte \texttt{\&dyn FeePolicy} fat pointers. You can now read a Rust type signature
and draw the exact stack and heap bytes it governs.

Yet engineers arriving from Java~6 and C often hit a second wall when writing their
first thousand lines of Rust. Their code compiles, yet every function starts with
four \texttt{let mut} accumulators and boolean flags, every conditional mutates an
outer variable across three branches, and every fallible lookup drowns in nested
\texttt{match} statements. They are writing imperative Java~6 in Rust syntax---and
the borrow checker pushes back because mutable state smeared across wide scopes
creates overlapping borrows.

Here is the historical connection that explains idiomatic Rust: although Rust took its
curly braces and flat memory model from C and C++, its control flow, algebraic types,
pattern matching, traits, closures, and error pipelines descend directly from two
functional languages---\textbf{OCaml} and \textbf{Haskell}. Read without the academic
jargon, both languages teach a small set of practical engineering habits:
compute by reducing expressions to values rather than mutating places step by step;
push side effects to explicit boundaries; make absence and failure ordinary values;
and compose small functions into pipelines. Studying OCaml and Haskell alongside Rust
shows where those ideas came from \emph{and} how ownership lets Rust compile them
into zero-allocation, register-tight machine loops without a garbage collector.

\chapterline

\section{Rewiring the \texttt{for}-Loop Brain: Statements vs.\ Expressions}
\label{ch07:statements-vs-expressions}

Java~6 and C split their grammar into two unequal categories: \emph{statements},
which execute in sequence to mutate memory places or jump elsewhere, and
\emph{expressions}, which evaluate to values (\texttt{a + b}, \texttt{x > 0}).
Because \texttt{if}/\texttt{else}, \texttt{switch}, \texttt{for}, and \texttt{\{...\}}
blocks are statements in Java~6, they cannot hand back values directly. That grammar
trains your fingers to compute by \emph{place-mutation}: declare an empty variable in
an outer scope, walk through branches that overwrite it, and hope every path left it
consistent.

Consider summarizing settled USD postings in a ledger batch using Chapter~6's
\texttt{FeePolicy} (\texttt{fee\_\allowbreak{}in\_\allowbreak{}cents}) to deduct fees on credits and assign a
reserve status string. In Java~6 we reach for outer mutable slots:

\begin{lstlisting}[language=Java]
// Java 6: computation via step-by-step place-mutation
public static BatchSummary summarizeUsd(
        List<Posting> postings, FeePolicy feePolicy) {
    long total = 0L;
    for (int i = 0; i < postings.size(); i++) {
        Posting p = postings.get(i);
        if (p.isSettled() && "USD".equals(p.getCurrency())) {
            if (p.isCredit()) {
                long gross = p.getAmountCents();
                total += gross - feePolicy.feeInCents(gross);
            } else {
                total -= p.getAmountCents();
            }
        }
    }

    String status;
    if (total >= 10000000L) {
        status = "INSTITUTIONAL_SURPLUS";
    } else if (total >= 0L) {
        status = "STANDARD_POSITIVE";
    } else {
        status = "DEFICIT_WATCH";
    }
    return new BatchSummary(total, status);
}
\end{lstlisting}

Notice three maintenance hazards hidden inside this everyday method:
\begin{enumerate}
  \item \textbf{Split declaration and initialization.} \texttt{String status;} creates
    a named slot before any value exists. If a teammate initializes
    \texttt{String status = null;} during a six-branch refactor to silence Java's
    definite-assignment check and forgets one \texttt{else if} arm, \texttt{null}
    escapes silently into production. In C or C++, reading an uninitialized
    \texttt{int64\_t total;} on a missed branch is immediate undefined behavior.
  \item \textbf{Temporal coupling and wide mutable scope.} Both \texttt{total} and
    \texttt{status} remain mutable until \texttt{return}. Reordering two statements
    below the loop can silently read half-updated state.
  \item \textbf{Entangled loop mechanics.} One \texttt{for} body mixes index stepping,
    USD settlement filtering, credit/debit sign logic, fee deduction, and summation.
\end{enumerate}

\subsection{Everything is an Expression: Java~6 vs.\ OCaml vs.\ Rust}

In OCaml and Rust, \texttt{if}/\texttt{else}, \texttt{match}, and \texttt{\{...\}}
blocks are \emph{expressions} that evaluate directly to typed values. Compare the
same settled-USD calculation in OCaml and in Rust:

\begin{lstlisting}[language=OCaml]
(* OCaml: match, if/else, and let-in blocks all evaluate to values *)
type direction = Credit | Debit
type posting = {
  settled : bool; usd : bool; dir : direction; cents : int64
}

let signed_net (fee_fn : int64 -> int64) (p : posting) : int64 =
  match p.dir with
  | Credit ->
      let fee = fee_fn p.cents in
      Int64.sub p.cents fee
  | Debit -> Int64.neg p.cents

let summarize_usd (posts : posting list) (fee_fn : int64 -> int64) =
  let total =
    posts
    |> List.filter (fun p -> p.settled && p.usd)
    |> List.map (signed_net fee_fn)
    |> List.fold_left Int64.add 0L
  in
  let status =
    if total >= 10000000L then "INSTITUTIONAL_SURPLUS"
    else if total >= 0L then "STANDARD_POSITIVE"
    else "DEFICIT_WATCH"
  in
  (total, status)
\end{lstlisting}

\begin{lstlisting}[language=Rust]
// Rust: blocks { ... tail_expr }, if, and match bind immutable values
#[derive(Clone, Copy, PartialEq, Eq)]
pub enum Direction { Credit, Debit }

#[derive(Clone, Copy, PartialEq, Eq)]
pub enum Currency { Usd, Eur }

pub struct Posting {
    pub settled: bool,
    pub currency: Currency,
    pub dir: Direction,
    pub amount_cents: i64,
}

pub trait FeePolicy {
    fn fee_in_cents(&self, amount_cents: i64) -> i64;
}

pub fn signed_net<P: FeePolicy>(policy: &P, p: &Posting) -> i64 {
    match p.dir {
        Direction::Credit => {
            let fee = policy.fee_in_cents(p.amount_cents);
            p.amount_cents - fee // tail expression: value of this block!
        }
        Direction::Debit => -p.amount_cents,
    }
}

pub fn summarize_usd<P: FeePolicy>(
    postings: &[Posting],
    policy: &P,
) -> (i64, &'static str) {
    let total: i64 = postings
        .iter()
        .filter(|p| p.settled && p.currency == Currency::Usd)
        .map(|p| signed_net(policy, p))
        .sum();

    let status = if total >= 10_000_000 {
        "INSTITUTIONAL_SURPLUS"
    } else if total >= 0 {
        "STANDARD_POSITIVE"
    } else {
        "DEFICIT_WATCH"
    };
    (total, status)
}
\end{lstlisting}

Every mutable outer slot disappeared. \texttt{total} and \texttt{status} are bound
once, immutably, at the exact instant their values are ready. Inside the
\texttt{Direction::Credit} arm, the block \texttt{\{ let fee = ...; p.amount\_cents - fee \}}
scopes \texttt{fee} locally, evaluates its tail expression without a semicolon, and
hands the resulting \texttt{i64} out as the value of the arm. Delete the final
\texttt{else} branch of \texttt{status}, and the program refuses to compile because
every arm of an expression must produce the exact same type.

\begin{focusbox}
\textbf{Statements Mutate Places; Expressions Reduce to Values.} Whenever you catch
yourself declaring \texttt{let mut result;} above an \texttt{if} or \texttt{match} in
Rust, rewrite it as \texttt{let result = match ... \{ ... \};}. You eliminate
uninitialized gaps, shrink mutable borrow lifetimes to zero, and force the type
checker to verify that every branch produces a compatible value.
\end{focusbox}

\subsection{Why \texttt{;} Yields \texttt{()} and Why C/Java \texttt{void} Breaks Generics}

Why does omitting the semicolon on \texttt{p.amount\_\allowbreak{}cents - fee} return
the \texttt{i64}, while adding \texttt{;} fails with ``expected \texttt{i64}, found
\texttt{()}''? Because Rust borrowed OCaml's \textbf{unit type} \texttt{()} instead of
C and Java's \texttt{void}.

In C and Java~6, \texttt{void} is a syntactic keyword meaning ``no value at all.'' You
cannot declare \texttt{void x;} or plug \texttt{void} into a type parameter \texttt{<T>}:
\begin{itemize}
  \item Java~6 had to invent \emph{both} \texttt{Runnable} (\texttt{void run()}) and
    \texttt{Callable<V>} (\texttt{V call()}), plus the uninstantiable class
    \texttt{java.lang.Void} whose sole value is \texttt{null}---forcing every
    \texttt{Callable<Void>} implementation to end with \texttt{return null;}.
  \item C++ libraries must specialize \texttt{Promise<void>} separately from
    \texttt{Promise<T>} because \texttt{T val = f();} fails to parse when \texttt{T}
    is \texttt{void}.
\end{itemize}

\begin{definition}[The Zero-Byte Unit Type \texttt{()}]
\label{ch07:def-unit-type}
The \emph{unit type} (written \texttt{unit} with sole value \texttt{()} in OCaml, and
\texttt{()} for both the type and its sole value in Haskell and Rust) is the 0-element
product type inhabited by exactly one value. Because distinguishing among $1$
possibility requires $\log_2(1) = 0$ bits, \texttt{std::mem::size\_of::<()>() == 0}.
\end{definition}

Making \texttt{()} a genuine zero-sized type (\emph{ZST}) unifies three rules cleanly:
\begin{enumerate}
  \item Writing \texttt{fn record\_metric(x: i64)} is exact shorthand for
    \texttt{fn record\_metric(x: i64) -> ()}.
  \item Placing \texttt{;} after an expression evaluates the expression, drops its
    result, and evaluates the block to \texttt{()}. Likewise, \texttt{if cond \{ work(); \}}
    without \texttt{else} compiles \emph{only} when its body yields \texttt{()},
    matching the implicit \texttt{else \{ () \}}.
  \item Generic types work uniformly with zero overhead: \texttt{Result<(), LedgerError>}
    carries zero payload bytes in its \texttt{Ok(())} variant, and Rust's standard
    \texttt{HashSet<K>} is a thin wrapper around \texttt{HashMap<K, ()>}
    generating zero value-storage bytes or instructions.
\end{enumerate}

\section{OCaml's Direct DNA in Rust}
\label{ch07:ocaml-dna}

\begin{historybox}
In the 1970s at Edinburgh, Robin Milner designed \textbf{ML} (\emph{Meta Language})
to script proof tactics inside the LCF theorem prover, inventing polymorphic
Hindley--Milner type inference and abstract module types so no buggy tactic could ever
forge an invalid theorem of type \texttt{thm}. At INRIA in France, Xavier Leroy,
Damien Doligez, J\'er\^ome Vouillon, and colleagues evolved Milner's ML through Caml into
\textbf{OCaml} (1996), pairing ML's expressive type system with a fast native-code
compiler (\texttt{ocamlopt}) and a low-latency generational garbage collector.

When \textbf{Graydon Hoare} started Rust in 2006, he wrote the entire original Rust
bootstrap compiler (\texttt{rustboot}, roughly 38{,}000 lines of code developed from
2006 to 2010 before \texttt{rustc} became self-hosting in 2011) \emph{in OCaml}. That
is why Rust's \texttt{let}, \texttt{fn}, \texttt{match}, \texttt{Option}/\texttt{Result},
\texttt{()} unit values, and tick-prefixed lifetimes (\texttt{'a}, echoing OCaml's
polymorphic type variables \texttt{'a}) carry unmistakable ML DNA.
\end{historybox}

Let us put real, runnable OCaml side by side with Java~6 and Rust to master five ML
ideas---and pinpoint exactly where Rust's lack of a garbage collector changes the
underlying data structures.

\subsection{Immutability by Default, Shadowing, and \texttt{ref} Cells vs.\ \texttt{let mut}}

In OCaml, \texttt{let} bindings are permanently immutable. Multi-step pipelines use
\emph{lexical shadowing} to bind clean intermediate stages, while genuine mutation
requires allocating an explicit mutable heap cell (\texttt{'a ref}) created with
\texttt{ref}, read with \texttt{!}, and overwritten with \texttt{:=}:

\begin{lstlisting}[language=OCaml]
(* OCaml: immutable shadowing vs explicit heap-allocated ref cell *)
let clamp_after_fee (raw_cents : int) (bps : int) : int =
  let cents = max 0 raw_cents in
  let fee   = (cents * bps) / 10000 in
  let cents = cents - fee in  (* shadows earlier cents; does not mutate! *)
  max 0 cents

let count_positive (xs : int list) : int =
  let count = ref 0 in (* allocates 16B mutable cell on GC heap *)
  List.iter (fun x -> if x > 0 then count := !count + 1) xs;
  !count
\end{lstlisting}

\begin{happeningbox}
\textbf{OCaml's 63-Bit Tagged Integers vs.\ Rust's Unboxed \texttt{let mut}.}
Because OCaml compiles polymorphic functions once without monomorphization, its
moving garbage collector must distinguish integers from pointers at runtime inside
any 64-bit register. OCaml reserves bit~0 of every 64-bit word as a tag: heap
pointers are 8-byte aligned (ending in \texttt{0}), while immediate \texttt{int}s,
\texttt{bool}s, and \texttt{()} set bit~0 to \texttt{1} (\texttt{(n << 1) | 1}),
giving 63-bit signed integers with zero boxing. (Full 64-bit \texttt{int64} values
box on the heap unless locally unboxed by \texttt{ocamlopt}.) Rust has no GC and
monomorphizes generics, so \texttt{let mut count: usize = 0;} mutates a raw 64-bit
stack/register integer in place with zero heap allocation, using borrow checking
(\texttt{\&mut count}) rather than a heap \texttt{ref} box when a closure mutates it.
\end{happeningbox}

\subsection{Algebraic Variants and Structural Pattern Matching}

Compare an OCaml algebraic variant representing credits, debits, and inline-record
reversals with its direct Rust \texttt{enum} mirror:

\begin{lstlisting}[language=OCaml]
(* OCaml: algebraic variant + structural match with 'when' guards *)
type posting =
  | Credit of int64
  | Debit of int64
  | Reversal of { orig_id : int; reason : string }

let net_delta (frozen : bool) (p : posting) : (int64, string) result =
  match (frozen, p) with
  | (_, Credit c) when c > 0L -> Ok c
  | (true, Debit _) -> Error "account frozen"
  | (false, Debit c) when c > 0L -> Ok (Int64.neg c)
  | (false, Reversal { orig_id; reason }) when orig_id > 0 ->
      if String.length reason >= 3 then Ok 0L else Error "short reason"
  | _ -> Error "invalid posting"
\end{lstlisting}

\begin{lstlisting}[language=Rust]
// Rust: identical variant modeling and 2D tuple + guard matching
pub enum PostingOp {
    Credit(i64),
    Debit(i64),
    Reversal { orig_id: u32, reason: String },
}

pub fn net_delta(
    frozen: bool, p: &PostingOp,
) -> Result<i64, &'static str> {
    match (frozen, p) {
        (_, PostingOp::Credit(c)) if *c > 0 => Ok(*c),
        (true, PostingOp::Debit(_)) => Err("account frozen"),
        (false, PostingOp::Debit(c)) if *c > 0 => Ok(-*c),
        (false, PostingOp::Reversal { orig_id, reason })
            if *orig_id > 0 =>
        {
            if reason.len() >= 3 { Ok(0) } else { Err("short reason") }
        }
        _ => Err("invalid posting"),
    }
}
\end{lstlisting}

Instead of Java~6 \texttt{instanceof} ladders or sprawling Visitor boilerplates, both
languages match the pair \texttt{(frozen, p)} as a flat 2D truth table and verify
exhaustiveness at compile time.

\subsection{Persistent Data Structures and Structural Sharing (\texttt{x :: xs})}

If OCaml lists are immutable, does prepending an element onto a 10{,}000-item list
copy all 10{,}000 items? Never. An OCaml \texttt{'a list} is either empty (\texttt{[]})
or an immutable heap \textbf{cons cell} \texttt{x :: xs} storing head \texttt{x} and
a pointer to tail \texttt{xs}:

\begin{lstlisting}[language=OCaml]
(* OCaml: O(1) persistent updates share the exact same immutable tail *)
let base = [20; 30; 40] (* 20 :: 30 :: 40 :: [] -> 3 cons cells *)
let a    = 10 :: base   (* 1 new cons cell pointing to base *)
let b    = 99 :: base   (* 1 new cons cell ALSO pointing to base *)
\end{lstlisting}

As Figure~\ref{ch07:fig-structural-sharing} illustrates, only \textbf{five} cons cells
exist in RAM. Because every cell is permanently immutable, \texttt{a} and \texttt{b}
safely alias \texttt{base} without defensive \texttt{new ArrayList<>(base)} copies.

\begin{figure}[H]
\centering
\begin{tikzpicture}[font=\small,>=Stealth,
  rootbox/.style={draw=sourceblue!85!black,fill=sourceblue!10,rounded corners=2pt,thick,minimum width=1.45cm,minimum height=0.58cm,font=\footnotesize\ttfamily},
  ptr/.style={->,thick,gray!80!black}]

  \node[font=\bfseries\scriptsize,text=sourceblue!85!black] at (0.85, 2.25) {BINDINGS};
  \node[rootbox] (ra) at (0.85, 1.45) {a};
  \node[rootbox] (rb) at (0.85, -0.05) {b};
  \node[rootbox] (rbase) at (0.85, -1.25) {base};

  \node[font=\bfseries\scriptsize,text=quotientorange!85!black] at (3.45, 2.25) {NEW ($\Oh(1)$)};
  \draw[quotientorange!85!black,fill=quotientorange!12,thick,rounded corners=2pt] (2.55, 1.1) rectangle (4.35, 1.8);
  \draw[quotientorange!85!black,thick] (3.55, 1.1) -- (3.55, 1.8);
  \node[font=\footnotesize\ttfamily\bfseries] at (3.05, 1.45) {10};
  \fill[gray!80!black] (3.95, 1.45) circle (2pt);

  \draw[quotientorange!85!black,fill=quotientorange!12,thick,rounded corners=2pt] (2.55, -0.4) rectangle (4.35, 0.3);
  \draw[quotientorange!85!black,thick] (3.55, -0.4) -- (3.55, 0.3);
  \node[font=\footnotesize\ttfamily\bfseries] at (3.05, -0.05) {99};
  \fill[gray!80!black] (3.95, -0.05) circle (2pt);

  \node[font=\bfseries\scriptsize,text=targetgreen!80!black] at (8.6, 2.25) {SHARED IMMUTABLE TAIL (\texttt{base}, SAFE ALIASING)};
  \draw[targetgreen!80!black,fill=targetgreen!9,thick,rounded corners=2pt] (5.45, 0.35) rectangle (7.25, 1.05);
  \draw[targetgreen!80!black,thick] (6.45, 0.35) -- (6.45, 1.05);
  \node[font=\footnotesize\ttfamily\bfseries] at (5.95, 0.7) {20};
  \fill[gray!80!black] (6.85, 0.7) circle (2pt);

  \draw[targetgreen!80!black,fill=targetgreen!9,thick,rounded corners=2pt] (8.05, 0.35) rectangle (9.85, 1.05);
  \draw[targetgreen!80!black,thick] (9.05, 0.35) -- (9.05, 1.05);
  \node[font=\footnotesize\ttfamily\bfseries] at (8.55, 0.7) {30};
  \fill[gray!80!black] (9.45, 0.7) circle (2pt);

  \draw[targetgreen!80!black,fill=targetgreen!9,thick,rounded corners=2pt] (10.65, 0.35) rectangle (12.65, 1.05);
  \draw[targetgreen!80!black,thick] (11.65, 0.35) -- (11.65, 1.05);
  \node[font=\footnotesize\ttfamily\bfseries] at (11.15, 0.7) {40};
  \node[font=\scriptsize\ttfamily,text=gray!75!black] at (12.15, 0.7) {[]};

  \draw[ptr] (ra.east) -- (2.55, 1.45);
  \draw[ptr] (rb.east) -- (2.55, -0.05);
  \draw[ptr,quotientorange!85!black] (3.95, 1.45) to[out=0,in=155] (5.45, 0.83);
  \draw[ptr,quotientorange!85!black] (3.95, -0.05) to[out=0,in=205] (5.45, 0.57);
  \draw[ptr,sourceblue!80!black] (rbase.east) to[out=0,in=235] (5.65, 0.35);
  \draw[ptr] (6.85, 0.7) -- (8.05, 0.7);
  \draw[ptr] (9.45, 0.7) -- (10.65, 0.7);

  \node[draw=structurepurple!75!black,fill=structurepurple!7,rounded corners=3pt,thick,
        text width=12.5cm,align=left,inner sep=5pt,font=\scriptsize] at (6.4, -2.95) {
    \textbf{\color{structurepurple!90!black}GC Nursery vs.\ Rust \texttt{Vec<T>}:}
    Each OCaml cons cell occupies 24B (8B GC header + 8B head + 8B tail pointer). OCaml's bump-pointer minor nursery allocates \texttt{10 :: base} in $\sim$2 instructions and collects dead cells in bulk. Because Rust has \textbf{no GC}, sharing tails requires \texttt{Rc<Node<T>>} (adding 16B ref-count headers + pointer-chasing cache misses), so idiomatic Rust favors flat contiguous \texttt{Vec<T>} buffers and slices (\texttt{\&[T]}) paired with lazy iterators.
  };
\end{tikzpicture}
\caption{Persistent structural sharing in OCaml: \texttt{a = 10 :: base} and \texttt{b = 99 :: base} allocate one cons cell each while sharing the immutable tail \texttt{[20; 30; 40]}.}
\label{ch07:fig-structural-sharing}
\end{figure}

\subsection{Currying, Partial Application, Pipes (\texttt{|>}), and Tail Recursion}

Every multi-argument OCaml function is \textbf{curried}: \texttt{apply\_bps} below has
type \texttt{int -> int -> int}, meaning \texttt{int -> (int -> int)}. Supplying just
the first argument (\texttt{apply\_bps 25}) returns a partially applied closure
\texttt{int -> int} ready to drop straight into a \texttt{|>} (\texttt{let (|>) x f = f x})
pipeline alongside \texttt{List.filter}, \texttt{List.map}, and \texttt{List.fold\_left}:

\begin{lstlisting}[language=OCaml]
(* OCaml: currying, partial application, |> pipe, and tail recursion *)
let apply_bps (bps : int) (cents : int) : int =
  cents - ((cents * bps) / 10000)

let net_after_25bps : int -> int = apply_bps 25

let settled_batch_total (amounts : int list) : int =
  amounts
  |> List.filter (fun c -> c > 0)
  |> List.map net_after_25bps
  |> List.fold_left ( + ) 0

(* Explicit-accumulator tail recursion compiles to a register jump *)
let sum_tail_recursive (xs : int list) : int =
  let rec loop acc remaining =
    match remaining with
    | [] -> acc
    | x :: rest -> loop (acc + x) rest (* tail call: reuses frame! *)
  in
  loop 0 xs
\end{lstlisting}

Look at \texttt{List.fold\_left} and \texttt{loop}: instead of mutating an outer variable,
the running state rides explicitly in \texttt{acc}. Because \texttt{loop (acc + x) rest}
is the final action of the function (\textbf{tail position}), \texttt{ocamlopt}'s
\textbf{Tail-Call Optimization (TCO)} updates the CPU registers holding \texttt{acc}
and \texttt{remaining} and jumps back to the top of \texttt{loop}---consuming $\Oh(1)$
stack space across ten million nodes. (Rust does not \emph{guarantee} tail calls,
partly because RAII \texttt{Drop} destructors must run after a callee returns, so a
call that looks like it is in tail position often is not. Idiomatic Rust uses a loop
or \texttt{Iterator::fold}, which compiles directly to a machine loop.)

\subsection{Hindley--Milner Inference, \texttt{fn} Boundaries, and Modules}

Write a polymorphic list transformer in OCaml with zero type annotations:

\begin{lstlisting}[language=OCaml]
let rec map_postings xs f =
  match xs with
  | [] -> []
  | x :: rest -> (f x) :: (map_postings rest f)
(* Inferred: val map_postings : 'a list -> ('a -> 'b) -> 'b list *)
\end{lstlisting}

Using \textbf{Hindley--Milner type inference} (Algorithm~W), OCaml assigns fresh type
variables to \texttt{xs} and \texttt{f}, collects equations from \texttt{x :: rest}
and \texttt{f x}, and solves (\emph{unifies}) them to deduce the most general type
\texttt{'a list -> ('a -> 'b) -> 'b list}.

\begin{whybox}
\textbf{Why Rust Stops Hindley--Milner at \texttt{fn} Boundaries.}
Inside a function body, Rust uses Hindley--Milner-style unification (\texttt{let mut v = Vec::new();}
is inferred from a later \texttt{v.push(5\_i64)}). Yet Rust mandates explicit parameter
and return types on every \texttt{fn} header for three engineering reasons:
\begin{enumerate}
  \item \textbf{Strictly local error messages:} a mistake inside function $B$ always
    fails inside $B$ against $B$'s signature rather than silently changing $B$'s
    inferred type and triggering a baffling unification clash three modules away.
  \item \textbf{SemVer contracts:} editing a function body cannot accidentally break
    callers' type or lifetime bounds.
  \item \textbf{Tractable traits + parallel checking:} trait-based method lookup,
    lifetime subtyping, and deref coercions would make whole-program inference
    fragile and slow; explicit \texttt{fn} headers let \texttt{rustc} check each
    function body on its own.
\end{enumerate}
\end{whybox}

Even in OCaml, production engineers seal module boundaries behind explicit signatures
(\texttt{module type} / \texttt{sig ... end}), hiding internal representations just as
Rust uses \texttt{mod} with private struct fields:

\begin{lstlisting}[language=OCaml]
module type FEE_POLICY = sig
  type t
  val make_bps     : int -> t option
  val fee_in_cents : t -> int64 -> int64
end

module TieredBps : FEE_POLICY = struct
  type t = int (* hidden outside TieredBps! *)
  let make_bps b = if b >= 0 && b <= 10000 then Some b else None
  let fee_in_cents bps cents =
    Int64.div (Int64.mul cents (Int64.of_int bps)) 10000L
end
\end{lstlisting}

\section{Haskell's Big Ideas Without the Ivory-Tower Fog}
\label{ch07:haskell-ideas}

Released in 1990 by Simon Peyton Jones, Philip Wadler, Paul Hudak, John Hughes, and
colleagues, \textbf{Haskell} pushed functional design to a clean extreme: pure
functions everywhere, lazy evaluation by default, and ad-hoc overloading through
typeclasses. Let us strip the academic mist off Haskell's four biggest ideas and see
how each one maps into everyday Rust.

\subsection{Referential Transparency vs.\ Rust's Controlled Local Mutation}

\begin{definition}[Referential Transparency]
\label{ch07:def-referential-transparency}
An expression $e$ is \emph{referentially transparent} if replacing $e$ anywhere with
its evaluated value leaves program behavior completely unchanged. A function is
\emph{pure} if its output depends solely on its explicit arguments and it produces
zero observable side effects (no mutation of external state, no I/O).
\end{definition}

Pure functions give you algebraic substitution, zero hidden temporal coupling, trivial
unit tests, and race-free parallelism. Compare sorting and deduplicating ledger
amounts in pure Haskell versus owned-value Rust:

\begin{lstlisting}[language=Haskell]
-- Haskell: pure transformation allocating new immutable list nodes
import Data.List (sort)

sortAndDedup :: [Int] -> [Int]
sortAndDedup xs = dedup (sort xs)
  where
    dedup [] = []
    dedup [y] = [y]
    dedup (y1 : y2 : rest)
      | y1 == y2  = dedup (y2 : rest)
      | otherwise = y1 : dedup (y2 : rest)
\end{lstlisting}

\begin{lstlisting}[language=Rust]
// Rust: in-place L1-cache mutation inside, 100-percent pure to callers!
pub fn sort_and_dedup(mut items: Vec<i64>) -> Vec<i64> {
    items.sort_unstable(); // mutates flat buffer in place, 0 heap allocs
    items.dedup();         // linear in-place compaction
    items
}
\end{lstlisting}

Consider \texttt{sort\_and\_dedup(mut items: Vec<i64>) -> Vec<i64>}. Inside its
body, it mutates array slots in L1 cache with zero garbage-collector churn. Yet to
\emph{every caller in the program}, \texttt{sort\_and\_dedup} is \textbf{100 percent
referentially transparent}. Because \texttt{items} is moved in by value (and even on
borrowed helpers, \texttt{\&mut} forbids any external alias from existing), no outside
observer can witness the in-place swaps. This is the style PL engineers call
\emph{functional in the large, imperative in the small}.

\subsection{Lazy Evaluation vs.\ Zero-Allocation Demand-Driven Iterators}

Haskell evaluates expressions \emph{lazily} (call-by-need): an unevaluated expression
is packaged as a heap \textbf{thunk} \texttt{[code\_ptr | env]} and forced only when
downstream pattern matching demands it. That lets Haskell compose pipelines over
infinite lists cleanly:

\begin{lstlisting}[language=Haskell]
-- Haskell: [1..] is infinite, yet 'take 5' pulls only until 5 evens arrive
firstFiveEvenSquares :: [Int]
firstFiveEvenSquares =
  take 5 (filter even (map (\x -> x * x) [1..]))
-- Evaluates to [4, 16, 36, 64, 100] and halts!
\end{lstlisting}

\begin{trapbox}
\textbf{Heap Thunks and Haskell's \texttt{foldl} Space Leak.}
Pervasive laziness has a sharp hardware cost: every suspended computation is a heap
object, and lazy accumulators silently build long chains of thunks. In Haskell,
\texttt{foldl (+) 0 [1..1000000]} does \emph{not} add integers as it walks the list;
it allocates one million nested heap thunks \texttt{(((0 + 1) + 2) + ...)}, blowing
memory when forced at the end unless you switch to strict \texttt{Data.List.foldl'}.
\end{trapbox}

Look at how Rust achieves the exact same pull-based infinite pipeline in a strict
language with \textbf{zero heap thunks}:

\begin{lstlisting}[language=Rust]
// Rust: pull-based laziness encoded as flat stack state machines
pub fn first_five_even_squares() -> Vec<i64> {
    (1_i64..)                      // 8-byte RangeFrom<i64> on stack
        .map(|x| x * x)            // wraps RangeFrom in stack Map struct
        .filter(|x| (x & 1) == 0)  // wraps Map in stack Filter struct
        .take(5)                   // pulls .next() until 5 items match
        .collect()                 // [4, 16, 36, 64, 100]
}
\end{lstlisting}

Every expression in Rust is strict, so an \texttt{i64} is always a raw 8-byte integer,
never a thunk. Laziness is opt-in through the \texttt{Iterator} trait, where each
stage advances on stack registers one element at a time when \texttt{.next(\&mut self)}
is pulled.

\subsection{Typeclasses (\texttt{Eq}, \texttt{Ord}, \texttt{Show}, \texttt{FeePolicy}) $\to$ Rust Traits}

In 1989, Philip Wadler and Stephen Blott introduced Haskell \textbf{typeclasses} to
support principled ad-hoc polymorphism (\texttt{Eq}, \texttt{Ord}, \texttt{Show})
without OOP class hierarchies. Put Haskell's \texttt{class}/\texttt{instance} right
next to Chapter~6's Rust \texttt{trait}/\texttt{impl}:

\begin{lstlisting}[language=Haskell]
-- Haskell: typeclass + separate instance + constrained generic function
class FeePolicy a where
  feeInCents :: a -> Int -> Int

data FlatBps = FlatBps Int deriving (Eq, Ord, Show)

instance FeePolicy FlatBps where
  feeInCents (FlatBps bps) cents = div (cents * bps) 10000

netCredit :: FeePolicy p => p -> Int -> Int
netCredit policy cents = cents - feeInCents policy cents
\end{lstlisting}

\begin{lstlisting}[language=Rust]
// Rust: direct descendant -- traits + derive + monomorphized bounds
#[derive(PartialEq, Eq, PartialOrd, Ord, Debug)]
pub struct FlatBps(pub i64);

impl FeePolicy for FlatBps {
    fn fee_in_cents(&self, cents: i64) -> i64 {
        (cents * self.0) / 10_000
    }
}

pub fn net_credit<P: FeePolicy>(policy: &P, cents: i64) -> i64 {
    cents - policy.fee_in_cents(cents)
}
\end{lstlisting}

Rust traits are Wadler--Blott typeclasses compiled via C++-style monomorphization (or
explicit \texttt{\&dyn FeePolicy} fat pointers). Unlike Java~6 interfaces, which attach
to heap objects via \texttt{this} headers, typeclasses and traits attach to \emph{types}
---allowing static factories like \texttt{Default::default() -> Self}, retroactive
implementations on existing types, and global \emph{coherence} (orphan rules).

\subsection{Functors (\texttt{.map()}) and Monads (\texttt{.and\_then()} and \texttt{?})}

Suppose validating a ledger debit runs four fallible steps in sequence:
\texttt{parse\_id -> find\_account -> active\_limit -> check\_headroom}, followed by a
pure cent-to-dollar scaling step. In Java~6 returning \texttt{null} on failure, four
dependent steps create a deeply indented \texttt{if (x != null)} pyramid:

\begin{lstlisting}[language=Java]
// Java 6: nested null-check pyramid
public static Long checkDebitDollars(
        LedgerBook book, String rawId, long debitCents) {
    Long id = parseId(rawId);
    if (id != null) {
        Account acct = book.findAccount(id);
        if (acct != null) {
            Long limit = activeLimit(acct);
            if (limit != null) {
                Long headroom = checkHeadroom(limit, debitCents);
                if (headroom != null) {
                    return headroom / 100L;
                }
            }
        }
    }
    return null;
}
\end{lstlisting}

Functional languages flatten this staircase with two exact abstractions:

\begin{enumerate}
  \item \textbf{Functor (\texttt{fmap} in Haskell, \texttt{.map()} in OCaml/Rust):}
    lifts an \emph{infallible} function $f : A \to B$ into a context $F\langle A\rangle \to F\langle B\rangle$
    (\texttt{Maybe a}/\texttt{Option<A>}, \texttt{Either e a}/\texttt{Result<A, E>}, or
    \texttt{Iterator}), obeying identity (\texttt{x.map(|a| a) == x}) and composition
    (\texttt{x.map(f).map(g) == x.map(|a| g(f(a)))}---the mathematical license for
    \emph{loop fusion}).
  \item \textbf{Why \texttt{.map()} alone nests wrappers:} \texttt{find\_account} has
    type \texttt{u64 -> Option<\&Account>}. Calling \texttt{parse\_id(raw).map(|id| find\_account(book, id))}
    wraps the inner \texttt{Option} inside a second \texttt{Some}, yielding
    \texttt{Option<Option<\&Account>>}.
  \item \textbf{Monad (\texttt{>>=} ``bind'' in Haskell, \texttt{Option.bind} in OCaml,
    \texttt{.and\_then()} in Rust):} takes $M\langle A\rangle$ and a context-producing
    step $A \to M\langle B\rangle$, runs the step if a value is present, and
    \emph{flattens one layer} so the output stays a flat $M\langle B\rangle$.
\end{enumerate}

Compare Haskell's \texttt{>>=} and \texttt{do}-notation directly beside Rust's
\texttt{.and\_then()} and \texttt{?} operator on the exact same four-step pipeline:

\begin{lstlisting}[language=Haskell]
-- Haskell: monadic bind (>>=) and equivalent do-notation
checkDebitBind :: LedgerBook -> String -> Int -> Maybe Int
checkDebitBind book raw debit =
  parseId raw
    >>= (\acctId   -> findAccount book acctId)
    >>= (\acct     -> activeLimit acct)
    >>= (\limit    -> checkHeadroom limit debit)
    >>= (\headroom -> Just (div headroom 100))

checkDebitDo :: LedgerBook -> String -> Int -> Maybe Int
checkDebitDo book raw debit = do
  acctId   <- parseId raw
  acct     <- findAccount book acctId
  limit    <- activeLimit acct
  headroom <- checkHeadroom limit debit
  Just (div headroom 100)
\end{lstlisting}

\begin{lstlisting}[language=Rust]
// Rust: .and_then() + .map() chain vs identical '?' operator
pub fn check_debit_chain(
    book: &LedgerBook, raw: &str, debit_cents: i64,
) -> Option<i64> {
    parse_id(raw)
        .and_then(|id| book.find_account(id))
        .and_then(|acct| active_limit(acct))
        .and_then(|limit| check_headroom(limit, debit_cents))
        .map(|headroom| headroom / 100)
}

pub fn check_debit(
    book: &LedgerBook, raw: &str, debit_cents: i64,
) -> Option<i64> {
    let id = parse_id(raw)?;
    let acct = book.find_account(id)?;
    let limit = active_limit(acct)?;
    let headroom = check_headroom(limit, debit_cents)?;
    Some(headroom / 100)
}
\end{lstlisting}

Look at \texttt{checkDebitDo} and \texttt{check\_debit}: a Monad is simply a
\textbf{programmable semicolon}. Between each line, \texttt{<-} in Haskell and
\texttt{?} in Rust inspect the \texttt{Maybe}/\texttt{Option} (or \texttt{Either}/\texttt{Result})
tag, unwrapping \texttt{Some}/\texttt{Ok} values along the top rail and diverting
\texttt{None}/\texttt{Err} onto a short-circuit bypass rail (Figure~\ref{ch07:fig-railway-monad}).

\begin{figure}[H]
\centering
\begin{tikzpicture}[font=\small,>=Stealth,
  swbox/.style={draw=sourceblue!85!black,fill=sourceblue!9,rounded corners=3pt,thick,align=center,minimum width=2.2cm,minimum height=0.92cm,font=\scriptsize},
  purebox/.style={draw=structurepurple!85!black,fill=structurepurple!9,rounded corners=3pt,thick,align=center,minimum width=2.05cm,minimum height=0.92cm,font=\scriptsize},
  okline/.style={->,very thick,targetgreen!80!black},
  errline/.style={->,very thick,trapred!80!black}]

  \node[font=\scriptsize\ttfamily\bfseries,text=gray!80!black] (inp) at (-0.65, 1.05) {raw};
  \node[swbox] (s1) at (1.35, 1.05) {\textbf{\texttt{parse\_id}}\\[-1pt]{\tiny \texttt{str -> Option<u64>}}};
  \node[swbox] (s2) at (4.15, 1.05) {\textbf{\texttt{find\_account}}\\[-1pt]{\tiny \texttt{.and\_then} / \texttt{?}}};
  \node[swbox] (s3) at (6.95, 1.05) {\textbf{\texttt{check\_headroom}}\\[-1pt]{\tiny \texttt{.and\_then} / \texttt{?}}};
  \node[purebox] (s4) at (9.75, 1.05) {\textbf{\texttt{|h| h / 100}}\\[-1pt]{\tiny Functor \texttt{.map()}}};

  \node[font=\scriptsize\ttfamily\bfseries,text=targetgreen!80!black] (okout) at (11.95, 1.05) {Some(d)};
  \node[font=\scriptsize\ttfamily\bfseries,text=trapred!80!black] (errout) at (11.95, -0.6) {None};

  \draw[okline] (inp.east) -- (s1.west);
  \draw[okline] (s1.east) -- node[above,font=\tiny\ttfamily,text=targetgreen!80!black] {id} (s2.west);
  \draw[okline] (s2.east) -- node[above,font=\tiny\ttfamily,text=targetgreen!80!black] {lim} (s3.west);
  \draw[okline] (s3.east) -- node[above,font=\tiny\ttfamily,text=targetgreen!80!black] {h} (s4.west);
  \draw[okline] (s4.east) -- (okout.west);

  \draw[trapred!80!black,very thick] (1.35, -0.6) -- (11.0, -0.6);
  \draw[errline] (11.0, -0.6) -- (errout.west);
  \draw[errline] (s1.south) -- node[right,font=\tiny\ttfamily,text=trapred!80!black] {None} (1.35, -0.6);
  \draw[errline] (s2.south) -- node[right,font=\tiny\ttfamily,text=trapred!80!black] {None} (4.15, -0.6);
  \draw[errline] (s3.south) -- node[right,font=\tiny\ttfamily,text=trapred!80!black] {None} (6.95, -0.6);

  \node[font=\tiny,text=structurepurple!85!black,align=center] at (9.75, 0.18) {(pure $A \to B$:\\[-1pt]no switch track)};
  \node[font=\tiny\bfseries,text=targetgreen!80!black] at (5.6, 1.85) {TOP RAIL (\texttt{Some(T)} / \texttt{Ok(T)}): unwrapped value passed into next stage};
  \node[font=\tiny\bfseries,text=trapred!80!black] at (5.6, -1.02) {BYPASS RAIL (\texttt{None} / \texttt{Err(E)}): short-circuits remaining pipeline};
\end{tikzpicture}
\caption{Railway-oriented pipeline: monadic bind (\texttt{.and\_then} / \texttt{?}) splices two-track switches without nesting \texttt{Option<Option<T>>}, whereas functor \texttt{.map} lifts a single-track function $A \to B$ onto the top rail.}
\label{ch07:fig-railway-monad}
\end{figure}


\section{Zero-Cost Closures in Rust: Anonymous Structs and \texttt{Fn}, \texttt{FnMut}, \texttt{FnOnce}}
\label{ch07:rust-closures}

Every higher-order pipeline we wrote in OCaml and Haskell passed closures---anonymous
functions that capture variables from their enclosing scope (\texttt{fun p -> p.cents >= min\_cents}).
Before we see how Rust compiles iterator chains into single machine loops, we have to
answer a fundamental hardware question: \emph{where do a closure's captured variables
live in memory, and how does the CPU jump to its code?}

In C, a plain function pointer \texttt{bool (*pred)(const Posting*)} is just an 8-byte
instruction address. It has nowhere to store a captured threshold \texttt{min\_cents},
which is why C APIs like \texttt{qsort\_r} or \texttt{pthread\_create} force you to
pack state into a separate struct and thread an untyped \texttt{void *ctx} pointer
by hand.

In Java~6 (anonymous inner classes like \texttt{new Predicate<Posting>() \{ ... \}}
capturing \texttt{final long minCents}) and in OCaml and Haskell, the compiler packs
captured variables into a \textbf{heap-allocated closure record}: a GC header, a
function or vtable pointer, and the captured fields. Passing a closure passes a
pointer to that heap block, and invoking it performs an \emph{indirect call} through
the code pointer---which allocates on the heap and blocks the compiler from inlining
tiny arithmetic bodies unless a JIT compiler later devirtualizes the call site.

Rust borrows C++11's unboxed lambda insight and strengthens it with ownership: in
Rust, \textbf{every closure expression generates its own unique, anonymous compiler
\texttt{struct} type} that stores its captured variables inline on the stack with zero
heap allocation and dispatches statically by default.

\subsection{Exact Compiler Desugaring of Three Ledger Closures}

To remove all mystery from Rust closures, let us trace three closures in our financial
ledger---one reading a threshold, one counting rejected items, and one flushing an
owned audit batch---and write down the exact \texttt{struct} and method signature
\texttt{rustc} synthesizes for each:

\begin{lstlisting}[language=Rust]
// Three Rust closures with three distinct capture & call modes:
#[derive(Clone)]
pub struct AuditEntry { pub post_id: u64, pub note: String }

pub fn demo_three_closures(
    min_cents: i64,
    mut rejected_count: usize,
    staged_audit: Vec<AuditEntry>,
) {
    // 1. Borrows &min_cents immutably -> implements Fn, FnMut, FnOnce
    let above_floor = |cents: i64| cents >= min_cents;

    // 2. Mutates &mut rejected_count -> implements FnMut, FnOnce (NOT Fn)
    let mut track_rejects = |cents: i64| -> bool {
        if cents < min_cents {
            rejected_count += 1;
            false
        } else {
            true
        }
    };

    // 3. Moves staged_audit out on call -> implements FnOnce ONLY
    let flush_audit = || -> usize {
        let consumed: Vec<AuditEntry> = staged_audit; // moves Vec out!
        consumed.len()
    };

    let _ = above_floor(500);
    let _ = track_rejects(100);
    let _ = flush_audit();
}
\end{lstlisting}

Under the hood, \texttt{rustc} inspects how each closure body uses each outer variable
(choosing the least restrictive capture mode: shared borrow \texttt{\&T} first,
mutable borrow \texttt{\&mut T} second, move \texttt{T} third) and synthesizes three
ordinary flat structs right in \texttt{demo\_three\_closures}'s stack frame:

\begin{lstlisting}[language=Rust]
// Conceptual desugaring generated by rustc (simplified trait signatures):
pub trait MyFnOnce<Arg> {
    type Output;
    fn call_once(self, arg: Arg) -> Self::Output;
}
pub trait MyFnMut<Arg>: MyFnOnce<Arg> {
    fn call_mut(&mut self, arg: Arg) -> Self::Output;
}
pub trait MyFn<Arg>: MyFnMut<Arg> {
    fn call(&self, arg: Arg) -> Self::Output;
}

// 1. Closure 'above_floor': 8 bytes on stack (1 shared reference)
struct __ClosureEnv1<'a> { min_cents: &'a i64 }

impl<'a> MyFn<i64> for __ClosureEnv1<'a> {
    fn call(&self, cents: i64) -> bool { cents >= *self.min_cents }
}
// (Plus MyFnMut and MyFnOnce impls delegating to self.call(cents)!)

// 2. Closure 'track_rejects': 16 bytes on stack (&i64 + &mut usize)
struct __ClosureEnv2<'a, 'b> {
    min_cents: &'a i64,
    rejected_count: &'b mut usize,
}

impl<'a, 'b> MyFnMut<i64> for __ClosureEnv2<'a, 'b> {
    fn call_mut(&mut self, cents: i64) -> bool {
        if cents < *self.min_cents {
            *self.rejected_count += 1;
            false
        } else {
            true
        }
    }
}
// Not MyFn(&self): mutating *self.rejected_count requires &mut self!

// 3. Closure 'flush_audit': 24 bytes on stack (ptr + len + cap of Vec)
struct __ClosureEnv3 { staged_audit: Vec<AuditEntry> }

impl MyFnOnce<()> for __ClosureEnv3 {
    type Output = usize;
    fn call_once(self, _arg: ()) -> usize {
        let consumed: Vec<AuditEntry> = self.staged_audit; // moves Vec!
        consumed.len()
    }
}
// Not MyFnMut(&mut self): a 2nd call would read a moved-from Vec!
\end{lstlisting}

Look at what happens to both memory and instruction dispatch:
\begin{itemize}
  \item \textbf{Zero heap allocation.} \texttt{\_\_ClosureEnv2} occupies 16 stack bytes.
    Capture-free closures (\texttt{|x: i64| x * 2}) occupy \textbf{0 bytes} and coerce
    to bare C-style function pointers \texttt{fn(i64) -> i64}.
  \item \textbf{Static dispatch that LLVM can inline.} Passing \texttt{track\_rejects} to
    a generic \texttt{retain} bounded by \texttt{F: FnMut(i64) -> bool} monomorphizes it for
    \texttt{F = \_\_ClosureEnv2}; LLVM sees the exact call target statically and inlines
    the body straight into the loop in optimized builds (while \texttt{\&dyn Fn(i64) -> bool} or
    \texttt{Box<dyn Fn(i64) -> bool>} remains available when runtime heterogeneity is needed).
\end{itemize}

\subsection{Why the Hierarchy Runs \texttt{Fn: FnMut: FnOnce}, and the \texttt{move} Trap}

At first glance, the subtrait bound \texttt{Fn: FnMut: FnOnce} looks backward to many
engineers: ``Wait---doesn't \texttt{Fn} allow \emph{more} ways to call a closure than
\texttt{FnOnce}? Why is \texttt{Fn} a subtrait of \texttt{FnOnce}?''

Read the hierarchy as \textbf{callee requirements vs.\ closure capabilities}
(Figure~\ref{ch07:fig-closures-layout-hierarchy}):
\begin{enumerate}
  \item \textbf{\texttt{FnOnce(self)} is the weakest demand a function can place on a
    callback:} ``Give me something I can call \emph{at most once} by consuming it''
    (used by \texttt{Option::unwrap\_or\_else} and \texttt{std::thread::spawn}).
    \emph{Every} closure in Rust satisfies \texttt{FnOnce}. Why? Because if you hand
    \texttt{unwrap\_or\_else} a pure \texttt{Fn(\&self)} closure like
    \texttt{above\_floor}, \texttt{unwrap\_or\_else} takes ownership of \texttt{self},
    borrows \texttt{\&self} once to run \texttt{Fn::call}, and drops it. Thus the set
    of \texttt{Fn} closures is a \emph{subset} of \texttt{FnOnce} closures ($\texttt{Fn} \subset \texttt{FnMut} \subset \texttt{FnOnce}$).
  \item \textbf{\texttt{FnMut(\&mut self)} asks for sequential repeatability:} ``I may
    call your closure zero, one, or many times in sequence, giving it unique
    \texttt{\&mut self} access each time'' (used by \texttt{Iterator::map},
    \texttt{Iterator::filter}, and \texttt{Vec::retain}). Both \texttt{FnMut} and
    \texttt{Fn} closures qualify; consuming \texttt{FnOnce}-only closures are rejected.
  \item \textbf{\texttt{Fn(\&self)} is the strictest demand:} ``I may call your closure
    many times through shared \texttt{\&self} borrows, possibly from several threads at
    once when the closure is also \texttt{Sync}'' (used by
    \texttt{rayon}'s parallel \texttt{par\_\allowbreak{}iter().\allowbreak{}map(...)}
    across CPU cores). Only closures that never mutate or consume captures qualify.
\end{enumerate}

\begin{trapbox}
\textbf{The Classic \texttt{move |...|} Myth: Capture Mode $\neq$ Call Trait.}
One of the most persistent misconceptions in Rust is believing that writing
\texttt{move |x| ...} turns a closure into an \texttt{FnOnce}-only closure. \textbf{False.}
\begin{itemize}
  \item The \texttt{move} keyword controls \emph{how captured variables enter the
    closure struct} when the closure is constructed (moving ownership into the struct's
    fields instead of storing \texttt{\&} borrows, so the closure can outlive the
    current stack frame).
  \item \texttt{Fn}, \texttt{FnMut}, and \texttt{FnOnce} are determined solely by
    \emph{what the closure body does to its fields when called}.
\end{itemize}
Look at this three-line proof:
\begin{lstlisting}[language=Rust]
let currency = String::from("USD");
// 'currency' is MOVED into the closure struct, yet body only reads &str!
let is_ccy = move |code: &str| code == currency.as_str();
assert!(is_ccy("USD") && !is_ccy("EUR")); // implements Fn -- call 1000x!
\end{lstlisting}
Here \texttt{is\_ccy} owns a 24-byte \texttt{String} field inside its struct (so it is
\texttt{'static} and safe to send to another thread), yet calling \texttt{is\_ccy("USD")}
only borrows \texttt{\&self.currency}. Therefore \texttt{is\_ccy} implements \textbf{\texttt{Fn}},
\texttt{FnMut}, \emph{and} \texttt{FnOnce}.
\end{trapbox}

\begin{figure}[H]
\centering
\begin{tikzpicture}[font=\small,>=Stealth]
  % LEFT PANEL: Memory Layout Comparison
  \node[font=\bfseries\scriptsize,text=sourceblue!90!black] at (2.8, 3.05) {MEMORY \& DISPATCH LAYOUT};

  % Java / OCaml Heap Closure
  \node[draw=quotientorange!85!black,fill=quotientorange!9,rounded corners=3pt,thick,
        minimum width=5.5cm,minimum height=1.7cm] at (2.8, 1.85) {};
  \node[font=\scriptsize\bfseries,text=quotientorange!90!black,anchor=west] at (0.18, 2.45) {Java 6 / OCaml: heap closure};
  \node[draw=quotientorange!80!black,fill=white,thick,font=\tiny\ttfamily,minimum width=1.15cm,minimum height=0.52cm] (jptr) at (0.85, 1.55) {clos\_ptr};
  \node[draw=quotientorange!85!black,fill=white,thick,font=\tiny\ttfamily,align=center,minimum width=3.1cm,minimum height=0.78cm] (jheap) at (3.75, 1.55)
    {GC hdr $|$ code\_ptr $|$ env\\[-1pt]{\tiny indirect jump + heap alloc}};
  \draw[->,thick,quotientorange!85!black] (jptr.east) -- (jheap.west);

  % Rust Stack Closure
  \node[draw=targetgreen!85!black,fill=targetgreen!9,rounded corners=3pt,thick,
        minimum width=5.5cm,minimum height=1.75cm] at (2.8, -0.2) {};
  \node[font=\scriptsize\bfseries,text=targetgreen!90!black,anchor=west] at (0.18, 0.42) {Rust: unboxed stack closure};
  \node[draw=targetgreen!85!black,fill=white,thick,font=\tiny\ttfamily,align=center,text width=4.7cm] at (2.8, -0.45)
    {\_\_ClosureEnv2 on the stack (16B):\\[-1pt]
     [\&min\_cents | \&mut rejected]\\[-1pt]
     {\tiny\rmfamily concrete type, direct call, LLVM can inline}};

  % RIGHT PANEL: FnOnce >= FnMut >= Fn Nested Capability Rings
  \node[font=\bfseries\scriptsize,text=structurepurple!90!black] at (9.65, 3.05) {CAPABILITY SETS (\texttt{Fn: FnMut: FnOnce})};

  % Outer FnOnce
  \draw[draw=trapred!80!black,fill=trapred!7,rounded corners=5pt,thick] (6.2, -1.08) rectangle (13.1, 2.7);
  \node[font=\scriptsize\bfseries,text=trapred!85!black,anchor=north west] at (6.35, 2.62)
    {\texttt{FnOnce(self)}: callable once};
  \node[font=\tiny\ttfamily,text=gray!85!black,anchor=north west] at (6.35, 2.22)
    {APIs: Option::unwrap\_or\_else, thread::spawn};

  % Middle FnMut
  \draw[draw=sourceblue!85!black,fill=sourceblue!8,rounded corners=4pt,thick] (6.55, -0.92) rectangle (12.75, 1.65);
  \node[font=\scriptsize\bfseries,text=sourceblue!90!black,anchor=north west] at (6.7, 1.57)
    {\texttt{FnMut(\&mut self)}: repeated calls};
  \node[font=\tiny\ttfamily,text=gray!85!black,anchor=north west] at (6.7, 1.18)
    {APIs: Iterator::map, filter, fold, Vec::retain};

  % Inner Fn
  \draw[draw=targetgreen!85!black,fill=targetgreen!12,rounded corners=3pt,thick] (6.9, -0.75) rectangle (12.4, 0.62);
  \node[font=\scriptsize\bfseries,text=targetgreen!90!black,anchor=north west] at (7.05, 0.54)
    {\texttt{Fn(\&self)}: shared calls};
  \node[font=\tiny\ttfamily,text=gray!85!black,anchor=north west] at (7.05, 0.14)
    {APIs: rayon::par\_iter, \&dyn Fn callbacks};
  \node[font=\tiny,text=targetgreen!85!black,anchor=north west] at (7.05, -0.25)
    {Every \texttt{Fn} closure also satisfies \texttt{FnMut} \& \texttt{FnOnce}!};
\end{tikzpicture}
\caption{Left: Heap-allocated closure indirection in Java~6/OCaml versus Rust's flat stack-allocated closure struct. Right: Nested closure trait sets mapping receiver borrows (\texttt{self}, \texttt{\&mut self}, \texttt{\&self}) to standard combinators.}
\label{ch07:fig-closures-layout-hierarchy}
\end{figure}

\section{Zero-Cost Iterator Pipelines, \texttt{Result} Transposition, and Loop Fusion}
\label{ch07:iterator-pipelines}

Now put stack-allocated closures together with pull-based iterators, and look at what
happens when Rust compiles a multi-stage functional pipeline.

In strict OCaml (\texttt{List.filter} followed by \texttt{List.map}) or JavaScript
(\texttt{arr.filter(...).map(...)}), each stage runs \emph{eagerly} from start to
finish, allocating a brand-new heap list or array of intermediate results before the
next stage even begins. On a batch of $100{,}000$ ledger postings, chaining two
\texttt{filter}s and two \texttt{map}s allocates four intermediate heap collections
and traverses RAM five times.

In Rust, the entire iterator ecosystem rests on one tiny trait with one required
state-machine method:

\begin{lstlisting}[language=Rust]
pub trait Iterator {
    type Item;
    fn next(&mut self) -> Option<Self::Item>;
    // 70+ default combinator methods (map, filter, fold, collect, ...)
}
\end{lstlisting}

Trace every type and byte when we run our settled-USD batch fold in Rust:

\begin{lstlisting}[language=Rust]
pub fn net_settled_usd<P: FeePolicy>(
    postings: &[Posting], policy: &P,
) -> i64 {
    postings
        .iter()
        .filter(|p| p.settled && p.currency == Currency::Usd)
        .map(|p| signed_net(policy, p))
        .fold(0_i64, |acc, delta| acc + delta)
}
\end{lstlisting}

What sits in memory right before \texttt{.fold(...)} runs?
\begin{enumerate}
  \item \texttt{postings.iter()} creates \texttt{std::slice::Iter<'\_, Posting>}---just
    two raw pointers (\texttt{ptr} and \texttt{end}, 16 bytes on the stack).
  \item \texttt{.filter(...)} wraps that \texttt{Iter} plus the 0-byte stateless
    predicate closure inside a stack struct \texttt{Filter<Iter<'\_, Posting>, Pred>}
    (still 16 bytes).
  \item \texttt{.map(...)} wraps \texttt{Filter} plus the 8-byte closure struct holding
    \texttt{\&policy} inside \texttt{Map<Filter<\allowbreak{}Iter<'\_, Posting>, Pred>, Mapper>}
    (24 bytes flat on the stack, \textbf{zero heap allocations}).
  \item Finally, \texttt{fold} runs \texttt{while let Some(d) = self.\allowbreak{}next() \{ acc += d; \}}.
    Because every adapter and closure type is statically known, LLVM inlines
    \texttt{Map::next}, \texttt{Filter::next}, and \texttt{Iter::next} into
    \textbf{one fused machine loop} (Figure~\ref{ch07:fig-eager-vs-fusion})---often
    faster than \texttt{for i in 0..\allowbreak{}postings.len()} because advancing
    \texttt{ptr} up to \texttt{end} eliminates per-step array bounds checks.
\end{enumerate}

\begin{figure}[H]
\centering
\begin{tikzpicture}[font=\small,>=Stealth,
  ebox/.style={draw=trapred!80!black,fill=trapred!8,rounded corners=2pt,thick,align=center,minimum height=0.68cm,font=\scriptsize},
  fbox/.style={draw=targetgreen!85!black,fill=targetgreen!10,rounded corners=2pt,thick,align=center,text width=12.35cm,inner sep=4pt,font=\scriptsize}]

  % Top: Eager Pipeline
  \node[font=\bfseries\scriptsize,text=trapred!85!black,anchor=west] at (0.0, 2.05)
    {Eager Collection Pipeline (OCaml \texttt{List.map} / JS \texttt{Array.map}): 2 Heap Allocs, 3 RAM Passes};
  \node[ebox,minimum width=1.9cm] (e0) at (1.0, 1.35) {Input List};
  \node[ebox,minimum width=1.9cm] (e1) at (3.7, 1.35) {\texttt{filter}};
  \node[ebox,minimum width=2.2cm,fill=quotientorange!15,draw=quotientorange!85!black] (e2) at (6.55, 1.35) {Heap Alloc \#1};
  \node[ebox,minimum width=1.7cm] (e3) at (9.15, 1.35) {\texttt{map}};
  \node[ebox,minimum width=2.2cm,fill=quotientorange!15,draw=quotientorange!85!black] (e4) at (11.75, 1.35) {Heap Alloc \#2};
  \draw[->,thick,trapred!80!black] (e0) -- (e1);
  \draw[->,thick,trapred!80!black] (e1) -- (e2);
  \draw[->,thick,trapred!80!black] (e2) -- (e3);
  \draw[->,thick,trapred!80!black] (e3) -- (e4);

  % Bottom: Rust Fused Pull Pipeline
  \node[font=\bfseries\scriptsize,text=targetgreen!85!black,anchor=west] at (0.0, 0.35)
    {Rust Fused Iterator State Machine (\texttt{Map<Filter<Iter>>}): 0 Heap Allocs, 1 Register Loop};
  \node[fbox] at (6.38, -0.35)
    {Contiguous \texttt{\&[Posting]} slice $\longrightarrow$ \textbf{Single LLVM-Fused Loop:} advance \texttt{ptr}, branch on \texttt{settled \&\& Usd}, subtract fee in register, add to \texttt{acc}};
\end{tikzpicture}
\caption{Eager pipelines materialize a full heap collection between every stage; Rust's pull-based adapter structs inline and fuse into a single allocation-free pass over flat memory.}
\label{ch07:fig-eager-vs-fusion}
\end{figure}

\subsection{Ownership Entry Points, Combinators, and \texttt{Result} Transposition}

To read and write idiomatic Rust pipelines effortlessly, keep three groups of tools at
your fingertips:

\begin{enumerate}
  \item \textbf{Three ownership entry points (mirroring \texttt{\&T}, \texttt{\&mut T}, \texttt{T}):}
    \begin{itemize}
      \item \texttt{xs.iter()} yields shared references \texttt{\&T}, leaving \texttt{xs} intact afterwards.
      \item \texttt{xs.iter\_mut()} yields exclusive references \texttt{\&mut T} to update elements in place.
      \item \texttt{xs.into\_iter()} consumes \texttt{xs} and moves owned \texttt{T} values out one by one.
    \end{itemize}
  \item \textbf{Work-saving structural adapters:}
    \begin{itemize}
      \item \texttt{filter\_map(|x| ... -> Option<B>)} fuses filtering and transforming
        into one pattern match---never write \texttt{.filter(|x| x.is\_some()).map(|x| x.unwrap())}.
      \item \texttt{flat\_map(|entry| entry.postings.iter())} flattens nested slices
        (one journal entry $\to$ its double-entry legs) without allocating intermediate \texttt{Vec}s.
      \item \texttt{enumerate()}, \texttt{zip(other)}, \texttt{chunks(64)}, and
        \texttt{windows(2)} eliminate manual index arithmetic (for instance,
        \texttt{balances.windows(2).map(|w| w[1] - w[0])} computes daily balance deltas
        with zero risk of \texttt{i - 1} underflow).
    \end{itemize}
  \item \textbf{Short-circuiting fallible iteration (\texttt{try\_fold} and \texttt{collect::<Result<Vec<T>, E>>()}):}
    What if parsing or arithmetic inside a pipeline can fail with \texttt{Result<T, E>}?
    Rust's \texttt{From\allowbreak{}Iterator} trait implements a clean \textbf{type transposition}:
    an \texttt{Iterator<Item = Result<T, E>>} can be collected directly into a
    \texttt{Result<Vec<T>, E>}.
\end{enumerate}

\begin{lstlisting}[language=Rust]
// Short-circuiting Result transposition + O(1)-space checked try_fold
pub fn parse_and_checked_sum(
    raw_lines: &[&str],
) -> Result<i64, &'static str> {
    // 1. Transposes Iterator<Item = Result<i64, E>> -> Result<Vec<i64>, E>
    //    Stops pulling IMMEDIATELY on the first Err!
    let amounts: Vec<i64> = raw_lines
        .iter()
        .map(|s| s.trim().parse::<i64>().map_err(|_| "bad integer"))
        .collect::<Result<Vec<i64>, &'static str>>()?;

    // 2. try_fold threads a fallible accumulator. On its own (fed by the
    //    parsing map directly) it needs no Vec at all -- see Exercise 7.1.
    amounts.iter().try_fold(0_i64, |acc, &x| {
        acc.checked_add(x).ok_or("i64 balance overflow")
    })
}
\end{lstlisting}

Look at \texttt{.collect::<Result<Vec<i64>, \&'static str>>()?}. If line~3 of
$10{,}000$ lines contains \texttt{"abc"}, \texttt{collect} sees \texttt{Err("bad integer")},
halts the iterator immediately without parsing lines~4 through $10{,}000$, drops the
partially populated \texttt{Vec}, and returns \texttt{Err}. And when you only need a
scalar aggregate rather than a \texttt{Vec}, \texttt{try\_fold} threads a fallible
accumulator in $\Oh(1)$ stack space with the exact same instant short-circuiting.

\subsection{Engineering Judgment: When a Plain \texttt{for} Loop Wins}

Thinking functionally does \emph{not} mean outlawing \texttt{for} loops or playing
combinator golf. Remember that \texttt{for x in xs} in Rust already desugars to
\texttt{IntoIterator::into\_iter(xs)}. Experienced Rust engineers prefer a plain
\texttt{for} loop over a contorted iterator chain in three clear situations:
\begin{enumerate}
  \item \textbf{Early function returns or multi-level loop exits:} if finding a fraud
    flag inside a loop should immediately \texttt{return Ok(Decision::Freeze)} from
    the enclosing function while updating a side log, \texttt{for} + \texttt{return}
    is far clearer than encoding exits inside \texttt{ControlFlow}.
  \item \textbf{Coupled multi-variable state machines:} when each step conditionally
    mutates three hash maps and a sliding watermark, threading a 4-tuple
    \texttt{(map1, map2, map3, mark)} through \texttt{.fold(...)} hurts readability.
    Encapsulate that local mutation inside a pure outer function (just like
    \texttt{sort\_and\_dedup}) and use a clean \texttt{for} loop inside.
  \item \textbf{Asynchronous \texttt{.await} points:} synchronous closures passed to
    \texttt{Iterator::map} cannot \texttt{.await} an SQLx query directly, whereas
    \texttt{for cmd in batch \{ execute(\&pool, cmd).await?; \}} interacts naturally
    with Chapter~12's async state machines and the \texttt{?} operator.
\end{enumerate}

\section{Exercises}
\label{ch07:exercises}

\begin{exercise}[Refactoring Imperative Java~6 into OCaml and Zero-Allocation Rust]
Consider a Java~6 method scanning raw cent strings: it skips blank strings, parses
integers, ignores non-positive numbers, deducts a 20-basis-point fee (\texttt{c - (c * 20) / 10000}),
and returns \texttt{null} if any non-blank string fails to parse or if the running
64-bit sum overflows.
\begin{enumerate}
  \item Write this pipeline in idiomatic \textbf{OCaml} using \texttt{int64},
    \texttt{List.filter}, and a recursive or \texttt{fold\_left} short-circuiting pass
    returning \texttt{int64 option}.
  \item Write the pipeline in \textbf{Rust} with \emph{zero} \texttt{Vec} allocations
    by combining \texttt{.iter()}, \texttt{.filter(...)} on trimmed slices, and
    \texttt{.try\_fold(0\_i64, ...)} returning \texttt{Option<i64>}. Verify that a
    malformed string at index~0 never parses index~1.
\end{enumerate}
\end{exercise}

\begin{exercise}[Counting Cons Cells and Pointer Overhead: OCaml vs.\ Rust]
Trace the following OCaml definitions on a 64-bit machine where each cons cell
occupies 24 bytes (8B GC header + 8B head + 8B tail pointer):
\begin{lstlisting}[language=OCaml]
let s0 = [100; 200; 300; 400]
let s1 = 50 :: s0
let s2 = 75 :: (List.tl s0)
\end{lstlisting}
\begin{enumerate}
  \item Draw the memory graph for \texttt{s0}, \texttt{s1}, and \texttt{s2}. How many
    cons cells and total heap bytes exist across all three lists combined?
  \item Suppose you model this exact shared-tail graph in safe Rust without a GC using
    \texttt{enum List \{ Nil, Cons(i64, Rc<List>) \}}. Given that \texttt{Rc<T>}'s
    heap allocation stores \texttt{(strong\_count: usize, weak\_count: usize, payload: T)},
    calculate the exact byte size of one \texttt{Rc<List>} node (accounting for
    \texttt{List}'s tag + alignment) and compare cache-line locality against a flat
    \texttt{Vec<i64>}.
\end{enumerate}
\end{exercise}

\begin{exercise}[Manual Rust Closure Desugaring and Trait Classification]
For each Rust closure below (assuming 64-bit Linux where \texttt{\&T} is 8B,
\texttt{String} and \texttt{Vec<T>} are 24B), write down (i) the exact fields of the
compiler-generated \texttt{struct \_\_ClosureEnv}, (ii) its stack byte size, and (iii)
which subset of \texttt{\{Fn, FnMut, FnOnce\}} it implements:
\begin{lstlisting}[language=Rust]
let rate: i64 = 35;
let c1 = |amt: i64| (amt * rate) / 10_000;

let allow: Vec<u64> = vec![1, 2, 9];
let c2 = move |id: u64| allow.contains(&id);

let mut max_seen: i64 = 0;
let c3 = |x: i64| { if x > max_seen { max_seen = x; } };

let logs: Vec<String> = vec![];
let c4 = || logs.into_iter().count();
\end{lstlisting}
Explain specifically why \texttt{c2} has \texttt{move} yet differs in trait capability
from \texttt{c4}.
\end{exercise}

\begin{exercise}[Custom Zero-Allocation Lazy \texttt{Iterator}]
Implement \texttt{Iterator<Item = i64>} for the following lazy adapter over any
signed-delta iterator \texttt{I: Iterator<Item = i64>}:
\begin{lstlisting}[language=Rust]
pub struct RunningBalances<I> {
    pub iter: I, pub balance: i64, pub tripped: bool,
}
\end{lstlisting}
On each \texttt{next} call, add the next delta to \texttt{balance}; if \texttt{balance < 0}
(an overdraft), set \texttt{tripped = true} and return \texttt{None} on that step and all
later calls; otherwise yield \texttt{Some(balance)}. Show that chaining this adapter
with \texttt{.take(3)} allocates zero heap bytes.
\end{exercise}

\begin{exercise}[Proving the Three Monad Laws on \texttt{Option::and\_then}]
For any pure Rust functions \texttt{f: fn(i64) -> Option<i64>} and
\texttt{g: fn(i64) -> Option<i64>}, state and prove by exhaustive case analysis on
\texttt{m = None} and \texttt{m = Some(v)} that \texttt{Option} obeys all three monad
laws:
\begin{enumerate}
  \item \textbf{Left identity:} \texttt{Some(x).and\_then(f) == f(x)}.
  \item \textbf{Right identity:} \texttt{m.and\_then(Some) == m}.
  \item \textbf{Associativity:} \texttt{m.and\_then(f).and\_then(g) == m.and\_then(|x| f(x).and\_then(g))}.
\end{enumerate}
Finally, give a concrete example with two impure closures \texttt{h1} and \texttt{h2}
that push a label into a shared \texttt{RefCell<Vec<String>>} log. Run them over a
three-element \texttt{Vec} first as two eager passes (map with \texttt{h1} and
\texttt{collect}, then map with \texttt{h2} and \texttt{collect}) and then as one
fused lazy chain, and show that the log records the side effects in a different
order. This is why rewriting pipelines freely (loop fusion) is only safe for pure
stages.
\end{exercise}
